% Generated from locale/tr/content/first-order-logic/axiomatic-deduction/deduction-theorem-quantifiers.tex; do not hand-edit.


\olfileid[tr]{fol}{axd}{ddq}

\olsection{Niceleyicilerle Tümdengelim Teoremi}

\begin{thm}[Tümdengelim Teoremi]
\ollabel{thm:deduction-thm-q} $\Gamma \cup \{!A\} \Proves !B$ ise
$\Gamma \Proves !A \lif !B$.
\end{thm}

\begin{proof}
Yine $\Gamma \cup \{!A\}$'dan $!B$ için verilen !!{derivation} ele alınır ve
uzunluğuna göre tümevarım yapılır.

Tümevarım başlangıcının ispatı, \olref[ded]{thm:deduction-thm} ispatındakiyle
aynıdır.

Tümevarım adımında, $\Gamma \cup \{!A\}$'dan $!B$ için bir !!{derivation}
ele alalım ve bunun bir çıkarım kuralıyla gerekçelendirilen $!B$ adımıyla
bittiğini varsayalım. Kural
modus ponens ise \olref[ded]{thm:deduction-thm} ispatındaki gibi ilerleriz.
Kural \QR{} ise $!B \ident !C \lif \lforall[x][!D(x)]$ olduğunu ve
$a$'nın $!C$, $!A$, $\Gamma$ ya da~$!D(x)$ içinde geçmediği durumda
$!C \lif !D(a)$ biçiminde !!a{formula} bulunduğunu ve bunun !!{derivation}
içinde daha önce geldiğini biliriz. Dolayısıyla
\begin{align*}
  \Gamma \cup \{!A\} & \Proves !C \lif !D(a),\\
  \intertext{ve tümevarım varsayımı uygulanabildiğinden}
    \Gamma & \Proves !A \lif (!C \lif !D(a)).\\
  \intertext{Ayrıca}
  & \Proves (!A \lif (!C \lif !D(a))) \lif ((!A \land !C) \lif !D(a))\\
  \intertext{olduğundan modus ponens ile}
  \Gamma & \Proves (!A \land !C) \lif !D(a).\\
  \intertext{Özdeğişken koşulu hâlâ geçerlidir; bu nedenle bu !!{derivation} içine \QR{} ile gerekçelendirilen bir adım ekleyerek}
    \Gamma & \Proves (!A \land !C) \lif \lforall[x][!D(x)].\\
  \intertext{Ayrıca}
    & \Proves ((!A \land !C) \lif \lforall[x][!D(x)]) \lif
    (!A \lif (!C \lif \lforall[x][!D(x)])),\\
  \intertext{dolayısıyla modus ponens ile}
    \Gamma & \Proves !A \lif (!C \lif \lforall[x][!D(x)])
\end{align*}
elde edilir; yani $\Gamma \Proves !A \lif !B$.

$!B$'nin \QR{} ile gerekçelendirildiği ancak
$\lexists[x][!D(x)] \lif !C$ biçiminde olduğu durumu alıştırma olarak
bırakıyoruz.
\end{proof}

\begin{prob}
  \olref[fol][axd][ddq]{thm:deduction-thm-q} ispatını tamamlayın.
\end{prob}

